% Cote 119: offprints, reports, courses, CV and letters (typed and printed part): archive cote 119, site page 194 (PDF page 195).
% typescript. Draft compared: none.
% AI-transcribed by Claude Opus 5.5 (claude-opus-5-5, adaptive thinking, effort high), Message Batches API, batch r119_full (msgbatch_01DwZNm5Y7YsKBWyAN4sFs8w), 2026-11-13; not yet checked by eye.
\scanpage{119/194}
\mspage{- 2 -}

conviction et lues sans plaisir ; un ou  deux laissés pour compte décidément
``irrécupérables'' - voilà à quoi se réduit presque le dérisoire bilan de
trois années d'activité enseignante.

Et maintenant ?

Que ferons nous, nous les nouveaux protagonistes, en cette nouvelle
année, hélas académique qui commence, pour répondre aux desiderata d'un cours
officiel, sans pour autant nous borner à reproduire le scénario immuable
du maître d'école pérorant devant ses écoliers ? Tout enseignement est cas-%
trateur, tout discours vain, qui ne s'adresse à des êtres dont la curiosité
déjà ne soit en éveil. Quand la curiosité est absente, et effacé peut-être
jusqu'au souvenir des temps reculés où elle était encore vivante en nous -
que faire alors pour lui redonner vie ? Ceci est notre première, notre prin-%
cipale question, celle qui doit précéder toute autre. Tant qu'elle reste en
suspens, tant que ne s'est éveillé en chacun le désir du Jeu - toute incita-%
tion à un voyage de découverte qui serait fait en commun reste entièrement
vide de sens.

Notre principal propos donc sera de provoquer au jeu l'enfant qui
sommeille dans l'Ecolier figé sur son bac, tout comme dans le Maître. Mais
est-ce bien au Maître d'école de provoquer ainsi - n'est-ce pas plutôt le rôle
de chacun de nous de provoquer tous les autres, à commencer par lui-même ?
Pour nous y inciter, ne faudrait-il, à défaut d'un intérêt  préalable pour une
``matière'' dont au fond un écolier n'a que f....., un sursaut de saine nausée
devant la perspective de reprendre encore et toujours le sempiternel ballet
mécanique, figurants falots dans le rite infiniment ressassé de notre propre
castration ! Ou alors, le rite aurait-il fini par agir, et bel et bien châtré
en nous l'homme et la femme libres et créateurs - serions-nous réduits vraiment
sans espoir au triste état d'Homonculus Studiensis ? En place alors, ``Maître''
et ``Ecoliers'', pour exécuter, soumis, votre tour de danse !

A nous de voir si nous serons l'enfant absorbé dans un jeu fasci-%
nant - ou sautillantes marionnettes...

\margin{194}
